% Abstract
% Author: Adnan Ghribi

The SPIRAL2 accelerator facility at GANIL relies on a complex LLRF system to maintain beam stability and quality. Operational challenges arise from multi-fault scenarios, requiring rapid and accurate fault diagnosis to minimize downtime. This paper presents a multi-phase machine learning pipeline for fault detection, root cause identification, and early fault-onset detection in LLRF systems. The pipeline employs a progressive architecture, starting with unsupervised anomaly detection (Isolation Forest, LOF, DBSCAN), followed by supervised multi-label classification, root cause analysis, and fault subtype discovery through clustering. Physics-informed feature engineering enhances model interpretability: a count-normalized SHapley Additive exPlanations (SHAP) analysis % MANUAL: "count-normalized" refers to controlling for the 20:1 physics-informed-vs-generic feature-count imbalance, detailed in Sec. 4/5, not restated here
 finds a physics-informed family dominant for four of seven fault categories -- e.g. effective cavity-decay features for cavity quench/breakdown (\MetricOneZeroPhysicsGroupShapCavityQuenchBreakdownTopGroupMeanPct\%) and modulator-command features for RF regulation out of tolerance (\MetricOneZeroPhysicsGroupShapRfRegulationOutOfToleranceTopGroupMeanPct\%) -- with generic statistical features dominant for the remaining three, including both categories with no dedicated physics feature family in this pipeline, pickup threshold (\MetricOneZeroPhysicsGroupShapPickupThresholdTopGroupMeanPct\%) and vacuum threshold (\MetricOneZeroPhysicsGroupShapVacuumThresholdTopGroupMeanPct\%). On \MetricZeroZeroDataAttritionChainVSixTotalEvents{} SPIRAL2 postmortem events, the pipeline achieves \MetricZeroEightPhaseThreeaRootcauseStabilityAccuracyMean\ $\pm$ \MetricZeroEightPhaseThreeaRootcauseStabilityAccuracyStd\ accuracy (repeated-split mean $\pm$ std) in root cause identification, \MetricZeroFivePhaseZeroPrecursorRandomForestRocAuc\ Receiver Operating Characteristic Area Under the Curve (ROC AUC) in detecting fault onset ahead of the interlock trip -- confidently detectable at least \MetricZeroFivePrecursorLeadTimeLeadTimeMsLowerBoundMeanCensored{}ms before the trip for sampled events, though the true detection horizon was not bounded by this measurement -- and discovers subtype structure (best $k=\MetricZeroNinecEnhancedSubclassDiscoveryRfAuthorizationAbsentBestK$ for every analyzed category) within \MetricZeroNinecEnhancedSubclassDiscoveryNCategoriesAnalyzed\ of \MetricZeroNinecEnhancedSubclassDiscoveryNCategoriesTotal\ fault categories with sufficient events to cluster. An explainability framework combining SHAP and Local Interpretable Model-Agnostic Explanations (LIME) provides physics-grounded interpretations. Measured single-event inference latency suggests feasibility for future real-time integration, pending further tail-latency and precursor-lead-time characterization; the modular architecture and the discipline of tracing every reported number to a reproducible source are, we argue, as transferable to other safety-critical accelerator systems as the specific results themselves.
